%% kapa: exam=1 round=32 year=2021 period=2 form=A subject=appraisal_law subject_ko=감정평가관계법규 seq=15 number=15
%% source: Q-Net 감정평가사 기출문제 (cst00309) artlSeq=5208397 file=extracted/1008/5208397/A형/2021년도 제32회 감정평가사 1차 2교시 A형.pdf
%% conversion: kordoc 4.24.51 -> tools/gen1.py (auto); review=proofread-sonnet; flags=none
\begin{question}{15}
부동산 가격공시에 관한 법령상 표준지공시지가에 관한 설명으로 옳지 않은 것은?
\begin{choices}
  \item 국토교통부장관은 표준지를 선정할 때에는 일단(一團)의 토지 중에서 해당 일단의 토지를 대표할 수 있는 필지의 토지를 선정하여야 한다.
  \item 국토교통부장관은 표준지공시지가를 공시하기 위하여 표준지의 가격을 조사·평가할 때에는 해당 토지 소유자의 의견을 들어야 한다.
  \item 국토교통부장관은 표준지공시지가의 조사·평가액 중 최고평가액이 최저평가액의 1.3배를 초과하는 경우에는 해당 감정평가법인등에게 조사·평가보고서를 시정하여 다시 제출하게 할 수 있다.
  \item 감정평가법인등은 표준지공시지가에 대하여 조사·평가보고서를 작성하는 경우에는 미리 해당 표준지를 관할하는 시·도지사 및 시장·군수·구청장의 의견을 들어야 한다.
  \item 표준지공시지가는 감정평가법인등이 제출한 조사·평가보고서에 따른 조사·평가액의 최저치를 기준으로 한다.
\end{choices}
\end{question}
\answer{5}
